\loai{Tính nhiệt độ}
\begin{vidu} %[1P5-3.2B][Loai1]
	Một bóng đèn dây tóc chứa khí trơ ở $27\doC$ và áp suất 0,6\ atm.\ Khi đèn sáng, áp suất không khí trong bóng là 1\ atm và không làm vỡ bóng đèn. Coi dung tích của bóng đèn không đổi, nhiệt độ của khí trong đèn khi cháy sáng là
	\choice
	{$380\doC$}
	{\True $227\doC$}
	{$450\doC$}
	{$500\doC$}
	
	\loigiai{
		Theo định luật Sác-lơ: 
		\begin{align*}
			\dfrac{0,6}{27 + 273} = \dfrac{1}{T_2} \Rightarrow T_2 = 500\ K = 227\doC.
		\end{align*}
	}
\end{vidu}
\loai{Tính áp suất}
\begin{vidu} %[1P5-3.2B][Loai1]
	Một bình thủy tinh kín chịu nhiệt chứa không khí ở điều kiện tiêu chuẩn. Nung nóng bình lên tới $200\doC$. Coi sự nở vì nhiệt của bình là không đáng kể. Cho $1\ atm = 1,013.10^5\ Pa$. Áp suất không khí trong bình là
	\choice
	{\True $1,755.10^5\ Pa$}
	{$1,6\ Pa$}
	{$1,633.10^5\ Pa$}
	{$1,755\ Pa$}
	\loigiai{
		\Binhluan{
			Điều kiện tiêu chuẩn: $t_1 = 0\doC,\ p_1 = 1\ atm = 1,013.10^5\ Pa.$
			\begin{align*}
				&\dfrac{p_1}{T_1} = \dfrac{p_2}{T_2}\\ 
				\Rightarrow\ &{p_2} = \dfrac{T_2}{T_1}{p_1} = \dfrac{200 + 273}{0 + 273}.1,013.10^5 = 1,755.10^5 \ Pa
			\end{align*}
		}
		{ĐKTC: $0^\circ C;\ 1\ atm$}
	}
\end{vidu}
\loai{Tính tỉ số}
\begin{vidu}%[1P5-3.2K][Loai3]
	Áp suất khí trơ trong bóng đèn tăng bao nhiêu lần khi đèn sáng nếu nhiệt độ đèn khi tắt là 25\doC, khi sáng là 323\doC.
	\choice
	{\True 2}
	{3}
	{4}
	{5}
	\loigiai{
		\Binhluan{
			\begin{itemize}
				\item Ta có: $T_1=298\ K;\ T_2=596\ K$
				\item Theo định luật Sác-lơ: $\dfrac{p_2}{p_1}=\dfrac{T_2}{T_1}=2.$
			\end{itemize}
		}
		{Kể cả khi tính tỉ số nếu không đổi đơn vị nhiệt độ ra độ K vẫn dẫn tới kết quả sai}
	}
\end{vidu}
\loai{Độ tăng áp suất}
\begin{vidu}%[1P5-3.2K][Loai4]
	Một bình kín chứa một lượng khí ở nhiệt độ $25\doC$ và áp suất 3 bar ($1\;bar = 10^5\;Pa$). Hỏi phải tăng nhiệt độ lên tới bao nhiêu độ để áp suất tăng gấp ba?
	\choice
	{890 K}
	{621 K}
	{$75\doC$}
	{\True $621\doC$}
	\loigiai{
		Theo định luật Sác-lơ: $\dfrac{p_1}{T_1}=\dfrac{p_2}{T_2}$.
		\begin{align*}
			\Rightarrow T_2=\dfrac{p_2(t_1+273)}{p_1}=\dfrac{9.10^5.(25+273)}{3.10^5}=894\;K=621\doC
		\end{align*}
	}
\end{vidu}
\loai{Độ tăng nhiệt độ}
\begin{vidu}%[1P5-3.2K][Loai5]
	Một bình kín thể tích không đổi chứa khí lí tưởng ở nhiệt độ 27\doC. Hỏi nhiệt độ trong bình tăng thêm một lượng là bao nhiêu? Biết áp suất ban đầu và sau khi nhiệt độ thay đổi lần lượt là 1 atm và 2,5 atm.
	\choice
	{ 320 K}
	{ \True 450 K}
	{ 350 K}
	{400 K }
	\loigiai{
		\Binhluan{
			\begin{itemize}
				\item Trạng thái 1: $T_1=300\ K;\ p_1=1\ atm$
				\item Trạng thái 2: $p_2=2,5\ atm$
				\item Theo định luật Sác-lơ: $\dfrac{p_1}{T_1}=\dfrac{p_2}{T_2} \Rightarrow T_2 = 750\ K$ 
				\item Vậy $\Delta T=T_2-T_1=450\ K$.
			\end{itemize}	
		}
		{Để tránh nhầm lẫn ta có thể ghi thông số p, T của từng trạng thái}
	}
\end{vidu}
\loai{Tăng thêm nhiệt độ - áp suất}
\begin{vidu}%[1P5-3.2K][Loai6]
	Khi đun nóng đẳng tích một khối khí thêm 1\doC thì áp suất khối khí tăng thêm $\dfrac{1}{360}$ áp suất ban đầu. Nhiệt độ ban đầu của khối khí đó là
	\choice
	{\True 87\doC}
	{360\doC}
	{350\doC}
	{361\doC}
	\loigiai{
		\Binhluan{
			\begin{itemize}
				\item Đề cho: $\left\{\begin{aligned}& \Delta p=\dfrac{1}{{360}}p_1 \\
					& \Delta T=T_2-T_1=t_2-t_1=1 
				\end{aligned}\right.$
				\item Áp dụng định luật Sác-lơ
				\begin{align*}
					\dfrac{{p_2}}{{p_1}}=\dfrac{{T_2}}{{T_1}}\Leftrightarrow \dfrac{{T_2-T_1}}{{T_1}}=\dfrac{{p_2-p_1}}{{p_1}} \\
					\begin{aligned}&{\Rightarrow T_1=\dfrac{{\Delta p}}{{\Delta T}}p_1=\dfrac{1}{{360p_1.1}}=360\;K} \\ &{\Rightarrow t_1=T_1-273={87}\doC} \\ \end{aligned}
				\end{align*}
			\end{itemize}
		}
		{Độ tăng nhiệt độ C $=$ độ tăng nhiệt độ K: $\Delta t=\Delta K$ .\\
			Trường hợp này nếu ta đổi $\Delta T=1+273=274\ K$ sẽ dẫn đến kết quả sai}
	}
\end{vidu}
\loai{Bài toán đồ thị}
\begin{vidu} %[1P5-3.2B][Loai7]
	immini[thm]{Một khối khí lí tưởng thực hiện hai quá trình như trên hình vẽ. Các thông số được cho trên đồ thị. Biết áp suất của chất khí khi bắt đầu quá trình là 12\ atm. Áp suất của khối khí khi kết thúc quá trình là
		
		\motcot
		{5\ atm}
		{2,667\ atm}
		{\True 13,33\ atm}
		{1,875\ atm}}{\begin{tikzpicture}[scale=1,thick,>=stealth']
			\tkzDefPoints{0/0/o, 4/0/t, 0/4/v, 3/3.2/1, 1.25/3.2/2, 1.25/1.2/3}
			\tkzDrawSegments[->](o,t o,v)
			\tkzDrawSegments[arrow data={0.5}{stealth}](1,2 2,3)
			\draw[dashed](1)--(3,0)node[below]{$600$};
			\draw[dashed](3)--(1.25,0)node[below]{$250$};
			\draw[dashed](3)--(0,1.2)node[left]{$2,4$};
			\draw[dashed](2)--(0,3.2)node[left]{$6,4$};
			\tkzLabelPoint[above](t){$T\ (K)$}
			\tkzLabelPoint[above](1){$(1)$}
			\tkzLabelPoint[above](2){$(2)$}
			\tkzLabelPoint[right](3){$(3)$}
			\tkzLabelPoint[right](v){$V\ (\ell)$}
			\tkzLabelPoint[below left](o){$O$}
	\end{tikzpicture}}
	\loigiai{
		\Binhluan{
			\begin{bang}[tabularx={l|X|X|X}]{}
				&\textbf{Trạng thái} $(1)$&\textbf{Trạng thái} $(2)$&\textbf{Trạng thái} $(3)$\\\hline
				\textbf{p}&$p_1=12\ atm$&$p_2$&$p_3$\\\hline
				\textbf{V}&$V_1=6,4\ l$&$V_2=6,4\ l$&$V_3=2,4\ l$\\\hline
				\textbf{T}&$T_1=600\ K$&$T_2=250\ K$&$T_3=250\ K$
			\end{bang}
		}
		{$(1)\to (2)$: Đẳng tích\\
			$(2)\to (3)$: Đẳng nhiệt}
		\begin{itemize}
			\item Từ $(1)$ đến $(2)$ là quá trình đẳng tích $\Rightarrow \dfrac{12}{600} = \dfrac{p_2}{250} \Rightarrow p_2 = 5\ atm$
			\item Từ $(2)$ đến $(3)$ là quá trình đẳng nhiệt $\Rightarrow 5.6,4 = {p_3}.2,4 \Rightarrow p_2 = 13,33\ atm$.
		\end{itemize}		
	}
\end{vidu}
\loai{Bài toán thực tế}
\begin{vidu} %[1P5-3.2T][Loai8]
	immini[thm]{Một bình đầy không khí ở điều kiện tiêu chuẩn ($0\doC$; $1,013.10^5\ Pa$) được đậy bằng một vật có khối lượng 2\ kg. Tiết diện của miệng bình là $10\ cm^2$ và gia tốc rơi tự do là $10\ m/s^2$. Tìm nhiệt độ lớn nhất của không khí trong bình để không khí không đẩy được nắp bình lên và thoát ra ngoài. Biết áp suất khí quyển là $p_0 = 10^5\ Pa$.
		\motcot
		{\True $50,4\doC$}
		{$323,4\doC$}
		{$121,3\doC$}
		{$115\doC$}
	}
	{
		\begin{tikzpicture}[scale=0.75,thick,>=stealth']
			\draw[fill=white](0,0)--++(0:2)--++(90:4)--++(180:2)--cycle;
			\foreach \i in {1,...,200}
			\fill[blue] (rnd*2cm, rnd*4cm) circle (0.8pt);
			\draw[purple,very thick](0,0)--++(0:2)--++(90:4)--++(180:2)--cycle;
			\draw[fill=gray,very thick](0,3.8)--++(0:2)--++(90:0.2)--++(180:2)--cycle;
		\end{tikzpicture}
	}
	\loigiai{
		\Binhluan{
			\begin{itemize}
				\item Lực tác dụng lên nắp bình gồm có áp lực của không khí trong bình ${\vec F _1}$, trọng lực $\vec P$, áp lực của khí quyển ${\vec F _0}$
				\item Để nắp bình không bị bật ra thì: 
				\begin{align*}
					P + {F_0} \ge {F_1}  \Leftrightarrow mg + {p_0}S \ge {p_2}S \Leftrightarrow \dfrac{mg}{S} + {p_0} \ge {p_2}\  (1)
				\end{align*}
				\item Khi nung bình thể tích không thay đổi, áp dụng định luật Sác-lơ: 
				\begin{align*}
					\dfrac{p_1}{T_1} = \dfrac{p_2}{T_2} \Rightarrow {p_2} = \dfrac{T_2}{T_1}{p_1}\ (2)
				\end{align*}
				\item Thay $(2)$ vào $(1)$: 
				\begin{align*}
					&\dfrac{mg}{S} + {p_0} \ge \dfrac{T_2}{T_1}{p_1}\\ \Rightarrow\ &{T_2} \le \left( {\dfrac{mg}{S} + {p_0}} \right)\dfrac{T_1}{p_1}\\ 
					= &\left( {\dfrac{2.10}{10.10^{- 4}} + 10^5} \right)\dfrac{273}{1,013.10^5} = 323,4\ K  = 50,4\doC.
				\end{align*}
			\end{itemize}
		}
		{Ta phân biệt \textbf{áp suất: p} và \textbf{áp lực: F}
			\begin{align*}
				F=pS
			\end{align*}
			Cũng có thể giải bằng cách hiểu áp suất bên ngoài gây ra lớn hơn áp suất bên trong bình
			\begin{align*}
				p_n \geq p_t \Leftrightarrow p_0+\dfrac{P}{S} \geq p_2 
			\end{align*}
		}
	}
\end{vidu}

